\documentclass[10pt,a4paper]{amsart}
\usepackage[margin=19mm]{geometry}
\usepackage{amsmath,amssymb,mathtools}
\usepackage[scr=rsfs,cal=euler]{mathalfa}
\usepackage{xfrac}
\usepackage[unicode,hidelinks,bookmarks=false]{hyperref}
\newcommand{\ie}{i.e.\ }
\newcommand{\cf}{cf.\ }
\newcommand{\resp}{respectively\ }
\newcommand{\Subsection}{\subsection}
\providecommand{\nobreakdash}{\nobreak\hbox{-}}
\setlength{\emergencystretch}{2em}
\multlinegap0pt
\usepackage{amssymb}
\usepackage{mathtools}
\usepackage{bm}
\usepackage{braket}
\usepackage[shortlabels]{enumitem}
\usepackage{xcolor}
\newtheorem{proposition}{Proposition}
\title{A Lie-Jordan Geometric Formulation of Lindblad Dynamics}
\author{Leonel Bixano, Victor Alberto Cruz-Barriguete, Guillermo L\'opez-Alvarez, V. G. Ibarra-Sierra, Jos\'e Luis Cardoso, Alejandro Kunold}
\date{Source: arXiv:2607.04052v1 (CC BY 4.0)}
\begin{document}
\maketitle

\noindent\emph{Source and licence.} A Lie-Jordan Geometric Formulation of Lindblad Dynamics. Version arXiv:2607.04052v1 (CC BY 4.0), distributed by arXiv under the Creative Commons Attribution 4.0 International licence, http://creativecommons.org/licenses/by/4.0/, as recorded in the arXiv licence metadata for this exact version and shown on the abstract page https://arxiv.org/abs/2607.04052v1. Full licence notice: \texttt{licenses/arxiv-2607.04052-CC-BY-4.0.txt}.\par\smallskip

\noindent\textit{Editorial scope.} Counted unit: the article's proposition eq:qubit-universal-D-blocks (source0.tex lines 1835--1875). The article's complete original proof of it is delivered with it (source0.tex lines 1877--1938, one \texttt{proof} environment of 62 lines). No mathematics is written, repaired or re-derived: the statement and the proof are the article's own lines, in the article's own notation, and no retained line is substituted or re-flowed.  Retained uncounted as the source-local context the proof needs: lines 1771--1803; lines 1772--1782.  Nothing was omitted from any retained span, so the package carries no omission record.  No retained line contains a citation, so the wrapper carries no bibliography and no bibliography entry.  No retained line contains a figure environment, an image-inclusion command or drawing code, so no image is delivered and the package declares no asset dependency.  Editorial formatting only, in addition: an AMS \texttt{amsart} preamble that restates the article's own declarations and macros used by the retained lines, namely the environments \texttt{proposition} on separate counters, together with the standard packages those lines need.  The retained environments print in this document as Proposition 1 in order of appearance, whereas the article prints the same items with the numbering of its own full document.  The complete native source of the article as distributed in the arXiv e-print for this exact version is bundled unchanged as \texttt{sources/204445/source0.tex}.
We choose the normalized Hermitian basis
\begin{equation}\label{eq:qubit-HS-basis}
    h_0
    =
    \frac{\mathbb I}{\sqrt2},
    \qquad
    h_a
    =
    \frac{\sigma_a}{\sqrt2},
    \qquad
    a=1,2,3,
\end{equation}
where \(\sigma_a\) are the Pauli matrices. This basis satisfies \(\operatorname{Tr}(h_\mu h_\nu) = \delta_{\mu\nu}\), using \(\sigma_a\sigma_b = \delta_{ab}\mathbb I + \mathrm{i}\epsilon_{ab}{}^c\sigma_c\), the ordered products become
\begin{subequations}\label{eq:qubit-ordered-products}
\begin{align}
    h_0h_0
    &=
    \frac1{\sqrt2}h_0,
    \\
    h_0h_a
    =
    h_ah_0
    &=
    \frac1{\sqrt2}h_a,
    \\
    h_ah_b
    &=
    \frac{\delta_{ab}}{\sqrt2}h_0
    +
    \frac{\mathrm{i}}{\sqrt2}
    \epsilon_{ab}{}^c h_c.
\end{align}
\end{subequations}

\begin{equation}
    h_0
    =
    \frac{\mathbb I}{\sqrt2},
    \qquad
    h_a
    =
    \frac{\sigma_a}{\sqrt2},
    \qquad
    a=1,2,3,
\end{equation}

\begin{proposition}[Universal qubit blocks]
For the basis \eqref{eq:qubit-HS-basis}, the universal elementary maps satisfy
\begin{subequations}\label{eq:qubit-universal-D-blocks}
\begin{align}
    \mathcal D_{00}(Z)
    &=
    0,
    \\
    \mathcal D_{a0}(Z)
    &=
    \frac1{2\sqrt2}
    [h_a,Z],
    \\
    \mathcal D_{0a}(Z)
    &=
    -\frac1{2\sqrt2}
    [h_a,Z].
\end{align}
\end{subequations}
For \(a,b,c=1,2,3\), their action on the basis elements is
\begin{subequations}\label{eq:qubit-Dab-action}
\begin{align}
    \mathcal D_{ab}(h_0)
    &=
    \mathrm{i}\epsilon_{ab}{}^k h_k,
    \label{eq:qubit-Dab-h0}
    \\
    \mathcal D_{ab}(h_c)
    &=
    \frac12
    \left(
    \delta_{ac}h_b
    +
    \delta_{bc}h_a
    \right)
    -
    \delta_{ab}h_c.
    \label{eq:qubit-Dab-hc}
\end{align}
\end{subequations}
\end{proposition}

\begin{proof}
Since \(h_0=\mathbb I/\sqrt2\),
\[
\mathcal D_{00}(Z)
=
\frac12Z
-
\frac12
\left\{
\frac{\mathbb I}{2},Z
\right\}
=
0.
\]
Likewise,
\begin{align*}
\mathcal D_{a0}(Z)
&=
\frac1{\sqrt2}h_aZ
-
\frac1{2\sqrt2}
\{h_a,Z\}
\\
&=
\frac1{2\sqrt2}
[h_a,Z],
\end{align*}
and the expression for \(\mathcal D_{0a}\) follows analogously.

For the remaining blocks,
\begin{align*}
\mathcal D_{ab}(h_0)
&=
h_ah_0h_b
-
\frac12\{h_bh_a,h_0\}
\\
&=
\frac1{\sqrt2}
(h_ah_b-h_bh_a)
\\
&=
\frac1{\sqrt2}[h_a,h_b]
\\
&=
\mathrm{i}\epsilon_{ab}{}^k h_k.
\end{align*}
Finally, repeated use of
Eq.~\eqref{eq:qubit-ordered-products} gives
\[
\mathcal D_{ab}(h_c)
=
\frac12
\left(
\delta_{ac}h_b
+
\delta_{bc}h_a
\right)
-
\delta_{ab}h_c.
\]
\end{proof}

\end{document}
